\documentclass[a4paper,10pt]{article}
\usepackage{amssymb}

%opening
\title{On Numerical Integration in CFS}
\author{Fabian Wein}

\begin{document}

\maketitle

\begin{abstract}
The numerical integration in CFS, based on the integration of the reference
elements, is a crucial factor for the
size of the system matrix. The numerical integration is also (improperly) blamed for
bad simulation results. This document summarizes some knowledge gained in reimplementing the relevant part in CFS and shows how to isolate the
effect of numerical integration on simulation.
\end{abstract}

\section{Mathematics}
Via numerical integration the integral over the test (Ansatz) functions is 
approximated. This is covered in any FEM book, advisable is \cite{Solin04}.  

\subsection{Economical Gauss Quadrature}
The authors in \cite{Solin04} claim to present the most efficient integration points available, called \textit{economical} integration. For CFS users this is called \textit{Gauss\_standard}. Other books give for rectangles (quadrilateral) and cubes (hexahedron, brick) the values
gained by applying the cartesian product of line integration values. Especially in 
$\mathbf{R}^3$ economical are much more efficient.
\subsection{The Right Order}
Obviously the order of integration depends on the order of the test functions. 
In CFS
all elements (LineFE, RectangleFE, ...) are derived from BaseFE and child classes depending on the order of the test function (Line1FE, Line2FE, Quad1FE, Quad2FE, ...). It is now
important to note, that it is not sufficient to always use integration rules of order 1 and 2! I found no explanation in the books \footnote{Please add it here if you know it in detail!} but taking the product of two bilinear functions the result has a higher order than one in the center of the element but practically vanishes in first order on the element boundary $\ldots$. The default values should be sufficient for most cases. Table \ref{tab:default_order} shows the current default values. In case you have difficulties with this defaults please start a discussion to alter this values!

\begin{table}
\begin{center}
\begin{tabular}{r|r|r|r|r|r}
element & test function & Gauss order & points \\
\hline
line & 1 & 1 & 1\\
     & 2 & 2=3 & 2 \\
triangle & 1 & 2 & 3 \\
         & 2 & 3 & 4 \\
rectangle & 1 & 2=3 & 4 \\
          & 2 & 2=3 & 4 \\
tetrahedron & 1 & 1 & 1 \\
         & 2 & 2 & 4 \\
hexahedron & 1 & 4=5 & 14 \\
         & 2 & 4=5 & 14 \\
\end{tabular}
\caption{This are the default integration orders when no other value is given in the
XML file. They should have sufficient accuracy for most cases. The test cases (section \ref{sec:test_cases}) show that the results don't get significantly better for higher orders but are void for lower orders. An exception is the triangle, where for one Gauss order less the results are also in the same order. Please note that these values are chosen only on the basis of three test cases!}
\label{tab:default_order}
\end{center}
\end{table}

\subsection{Higher Orders}
On bad simulation results, the impact of the numerical integration can easily be eliminated
by manually applying an extreme integration order. Table \ref{tab:order_points} shows the number of integration points for a first set of higher orders. Not all data from \cite{Solin04} is implemented, but as the book is accompanied by a CD further points can be added to the code quite easily. For elements where the integration rules are the same for even and odd order (2,3) (4,5), ... only the higher order is stored. So users should ask for it otherwise a warning message will be printed. 

The XML manual shows how to manually adjust for
high orders. This should eliminate the effect of numerical integration.

\begin{table}
\begin{center}
\begin{tabular}{r|r|r|r|r|r}
order & line & trian & quad & tetra & hexa \\
\hline
1     & 1    & 1     & 1    & 1     & 1    \\ 
2     & 2    & 3     & 4    & 4     & 6    \\
3     & 2    & 4     & 4    & 5     & 6    \\
4     & 3    & 6     & 8    & 11    & 14   \\
5     & 3    & 7     & 8    & 14    & 14   \\
6     & 4    & 12    & 12   & 24    & 27   \\
7     & 4    & 13    & 12   & 31    & 27   \\
8     & 5    & 16    & 20   & 43    & 53   \\
9     & 5    & 19    & 20   & 53    & 53   \\
\ldots & \ldots & \ldots & \ldots & \ldots & \ldots \\
\end{tabular}
\caption{Economical Gauss quadrature points for the main elements. Further orders are implemented and not quoted here. In \cite{Solin04} orders up to greater 20 are given. Note that for line, quad (rectangle) and hexa (cube) the odd and even orders coincide.}
\label{tab:order_points}
\end{center}
\end{table}

\subsection{Reduced Integration}
This is a very special numerical trick to avoid blocking effects. It is mentioned in \cite{Hughes00}. The parameters are set directly in the code. Anybody using this reduced integration should check and maintain this values!

It should be noted, that in the test cases a lower integration order than the one used as defaults produced completely wrong results or even a crash in the solver.

\section{Validation}
The numerical integration was tested on three examples where the analytical solution is known. 
\subsection{The Test Cases}
\label{sec:test_cases}
The test cases are taken from the \textit{ANSYS Verification Manual 5.3} and are part of the test suite. 
\subsubsection*{2D Example: VM181 (Plate)}
VM181 is a test for the natural frequency of a flat circular plate with a clamped edge. The plate has radius a=0.4318m and height of t=0.0127m. The material properties are E=2.068428e11 Pa, $\nu$=0.3 and $\rho$=7801.51. The exact result of the first mode of vibration is 172.64Hz.

The examples are meshed via GID structured (coarse and fine) and unstructured for triangles and quadrilaterals for bilinear and biquadratic test functions.

\subsubsection*{3D Example: VM184 (Straight Cantilever Beam)}
The cantilever has the geometry l=0.1524m, h=0.00508m, t=0.00254m and is fixed on one and. On the other side a force can be applied. In the line
\begin{verbatim}
<load name="load" dof="uz" value="4.44822/8"/>
\end{verbatim}
the direction of the force is to be set. As with GID one cannot apply the force on the surface, the total force is to be divided by the number of points on the surface. To find it,
simply calculate the problem and see how many points are actually written to the history directory - then adjust the divider in the XML file. The analytical solution of the deflection for a force of 4.44822 Newton in x-direction is 7.62e-7m, for y-direction 2.743e-3m and for z-direction 1.097e-2m. 

The test cases are modelled with 10 and 30 tetrahedrons and hexahedrons with test functions of first and second order.

\subsubsection*{3D Example: VM187 (Bending of a Curved Beam)}
The beam is here a quarter of a ring with inner radius 0.104684m and outer radius 0.109728m. The depth is 0.00254m and the force 4.44822 Newton. The analytical solution is 2.24927e-3m. As for the straight beam, check the division of the force in the XML file.

The meshes consist of 20 and 40 elements.

\subsection{Observations}
The observations clearly show that the results are determined by
\begin{itemize}
\item the order of the test function.
\item the discretization.
\item the quality of the meshing (the deformation of elements).
\item a sufficient high interpolation order.
\end{itemize}
This is neither a surprise nor the observation that higher integration orders have 
only a negligible effect due to their theoretical exact approximation of corresponding 	polynomials.

One exception is the simulation of the heavily deformed meshed 2D plate with quadrilaterals. Here the results become worse by $\approx 5 \%$ for the next integration order but improve with successive higher orders. But such a  deformed meshing should be avoided anyway. Similar meshing with triangles gives better results but again the asymptotical convergence is about 40\% away from the analytical solution. 

\begin{thebibliography}{widest-label}
\bibitem[Solin04]{Solin04} \emph{Higher-Order Finite Element Methods},
  Pavel Solin, Karel Segeth, Ivo Dolezel, 2004

\bibitem[Hughes00]{Hughes00} \emph{The Finite Element Method},
  Thomas J. R. Hughes, 2000

\end{thebibliography}

\end{document}
